% Full title: Literature Review and Novelty Audit for Periodic Line-Defect Waveguide Guided Modes
\documentclass[UTF8,11pt,fontset=none]{ctexart}

\setCJKmainfont[BoldFont={Songti SC Bold}]{Songti SC}
\setCJKsansfont{PingFang SC}
\setCJKmonofont{PingFang SC}

\usepackage[a4paper,margin=2.15cm]{geometry}
\usepackage{amsmath,amssymb,mathtools}
\usepackage{booktabs,longtable,tabularx,array,multirow}
\usepackage[table]{xcolor}
\usepackage{enumitem}
\usepackage{microtype}
\usepackage{xurl}
\usepackage{hyperref}
\usepackage{fancyhdr}
\usepackage{caption}
\usepackage{pdflscape}
\usepackage[backend=biber,style=numeric,sorting=nyt,maxbibnames=10,giveninits=true]{biblatex}

\addbibresource{references.bib}
\hypersetup{
  colorlinks=true,
  linkcolor=blue!55!black,
  citecolor=green!40!black,
  urlcolor=blue!60!black,
  pdftitle={Periodic line-defect waveguide 文献调研与创新性审计},
  pdfauthor={Codex literature audit}
}
\setlength{\parindent}{2em}
\setlength{\parskip}{0.25em}
\setlist[itemize]{leftmargin=2em,itemsep=0.25em,topsep=0.3em}
\setlist[enumerate]{leftmargin=2.2em,itemsep=0.25em,topsep=0.3em}
\renewcommand{\arraystretch}{1.18}
\newcolumntype{Y}{>{\raggedright\arraybackslash}X}
\newcolumntype{C}[1]{>{\centering\arraybackslash}p{#1}}
\newcommand{\statusA}{\textbf{A（已有）}}
\newcommand{\statusB}{\textbf{B（直接扩展）}}
\newcommand{\statusC}{\textbf{C（部分新）}}
\newcommand{\statusD}{\textbf{D（可能新但证据不足）}}
\newcommand{\draft}{\texttt{draft/}\allowbreak\texttt{draft.tex}}
\setlength{\headheight}{14pt}
\setlength{\emergencystretch}{1.5em}
\pagestyle{fancy}
\fancyhf{}
\fancyhead[L]{Periodic line-defect waveguide 文献与创新性审计}
\fancyhead[R]{检索截止：2026-07-20}
\fancyfoot[C]{\thepage}
\setcounter{tocdepth}{2}

\title{\bfseries Periodic Line-Defect Waveguide Guided Modes\\
系统文献调研与当前 Draft 创新性审计}
\author{针对 \texttt{tep} 项目的证据型审计报告}
\date{检索截止日期：2026 年 7 月 20 日}

\begin{document}
\maketitle

\begin{abstract}
本报告审计当前 \draft{} 中 periodic line-defect waveguide guided-mode/eigenvalue
问题的理论、边界积分、Bloch 模态和数值特征值构造。证据集包括本地
\texttt{ref/} 的 13 个 PDF（其中 12 篇外部学术论文、1 份内部旧报告）以及
13 篇经期刊页、DOI 或 arXiv 核实的外部论文；20 篇最直接相关论文另有逐篇笔记。
结论是：全局导模问题到小邻域非线性 DtN 问题的精确等价、吸收传播算子/Riccati
构造、RtR 对 forbidden frequencies 的处理、广义 Floquet/Jordan 模态完备性、
以及 smallest-singular-value 扫描均已有明确先例。draft 最可能保留的研究价值，
不是这些组成部分本身，而是“高阶 Müller 界面积分离散 + 广义稳定 Bloch Cauchy
relation + 可排除 representation nullspace 的特征值认证”这一特定交叉耦合。
然而当前 Theorem 2 的 Appendix A 在一个关键式子中混用了 (k) 与 (nk) 的
quasiperiodic double-layer operator，Conjecture 1 又遗漏 Jordan chains，因此尚不能把
该耦合宣称为已证明的新方法。报告给出六条候选方向、优先级、主路线、备用路线和明确
停止/转向条件。
\end{abstract}

\tableofcontents
\clearpage

\section{调研目标、范围与结论先行}

本次任务回答五个问题：哪些内容已经解决，draft 与哪些工作重合，哪些只是直接改写，
哪些交叉方向仍可能成文，以及下一步应优先投入什么。审计对象是 2026-07-20 21:04:36
修改的 \draft{} 及其 \texttt{appendixA.tex}；它比 \texttt{pre/} 下各候选版本更新，且
同时包含 outgoing Cauchy trace subspaces、Theorems 1--2、Conjectures 1--2、MFS
quasiperiodic Green function、Müller BIE、one-cell scattering matrix 和 Bloch pencil，
故不是仅凭文件名选择。

本报告使用如下等级：A 为已有结果，B 为已有方法的直接扩展，C 为部分新颖但需要证明，
D 为截至本次检索可能新但证据不足，E 为有充分证据的真正未解问题。这里未给任何现有
draft claim 评为 E，因为“没有找到直接重合”不能证明“此前无人研究”。

\begin{table}[htbp]
\centering
\caption{核心结论摘要}
\small
\begin{tabularx}{\textwidth}{p{0.24\textwidth} C{2.0cm} Y}
\toprule
Draft 组成 & 状态 & 审计结论 \\
\midrule
全局问题与 bounded trace/DtN 问题等价 & A & Fliss 2013 已证明 line-defect guided-mode 的精确小区域 DtN 等价且保持重数。 \\
完整 Cauchy relation & B/C & 正则处只是 DtN graph 的关系写法；exceptional frequencies 处可能有价值，但 RtR 已解决 forbidden-frequency 可计算性。 \\
Theorem 2 exact/unique representation & C（证明有误） & 证明框架与 Barnett--Greengard Appendix A 高度重合；当前混用 \(D^{(nk)}_{\rm QP}\) 与 \(D^{(k)}_{\rm QP}\)。 \\
Bloch trace completeness & A（需改写） & 广义 Floquet modes/Jordan chains 的 Riesz basis 已有定理；普通 eigenvectors 一般不够。 \\
MFS QP Green function 与三单元近场 & B & proxy、近邻自由空间源和 Rayleigh matching 均有先例；除非补稳定性/误差结果。 \\
scattering-to-Bloch pencil & A/B & 是标准 scattering/transfer/wave-FEM 代数构造；直接稳定 relation coupling 仍可研究。 \\
smallest singular value scan & A & 旧有启发式；应由 contour/Riesz projection 加物理残差认证替代。 \\
当前数值证据 & 预备性 & line-defect 小节为空，尚无完整耦合算例、收敛表或独立 solver 比较。 \\
\bottomrule
\end{tabularx}
\end{table}

\section{当前 draft 的问题与方法概述}

轴向 Bloch 参数固定为 \(\beta\)，横向无界的 line-defect eigenproblem 可抽象写为
\begin{equation}
  \mathcal L_\beta u = k^2 u,\qquad
  u\in H^1_{\beta}(\mathbb R\times(0,L)),\qquad
  u(x,\cdot)\to0\quad (|x|\to\infty),
\end{equation}
其中外侧周期 leads 与中间 defect cell 的系数不同。draft 的第一步是在左右 artificial
ports \(\Gamma_L,\Gamma_R\) 定义 outgoing Cauchy trace relations
\begin{equation}
  \mathcal C_\pm(k,\beta)
  \subset H^{1/2}(\Gamma_\pm)\times H^{-1/2}(\Gamma_\pm),
\end{equation}
再把全局问题化为 center cell 内 PDE 加
\((\gamma u,\partial_n u)|_{\Gamma_\pm}\in\mathcal C_\pm\)。若 Dirichlet 投影可逆，
\(\mathcal C_\pm=\{(f,\Lambda_\pm f)\}\) 是 DtN graph。

第二步以 quasiperiodic/free-space layer potentials 与独立 homogeneous Rayleigh field
表示单元场，利用 transmission jumps 形成 Müller-type block。第三步从 one-cell
scattering matrix \(S\) 和端口相位条件组成关于 Bloch multiplier \(\lambda\) 的 pencil，
选取 \(|\lambda|<1\) 的稳定 modes，再把其 Cauchy traces 作为半导外边界 relation 的
有限维基。最终非线性矩阵可示意为
\begin{equation}
  \mathcal A(k,\beta)
  \begin{bmatrix}\eta\\a_{\rm Rayleigh}\\c_L\\c_R\end{bmatrix}=0,
  \qquad
  \sigma_{\min}(\mathcal A(k,\beta))\approx0.
\end{equation}
现有数值内容验证 QP Green function 的 MFS 近似，并测试 homogeneous-lead special
case；尚未展示一般 line defect。

\section{检索方法、数据库与证据分层}

检索截止日期为 2026-07-20。关键词被分为七组：line-defect guided modes；DtN/RtR/
propagation operator；boundary integral；Bloch completeness/operator pencil；TE 与
discontinuous coefficients；band edge/embedded/leaky/Wood；nonlinear eigenvalue 与
contour methods。所有实际 query 和舍弃理由见伴随文件 \texttt{search\_log.md}。

Backward chaining 从 Joly--Li--Fliss、Fliss、Barnett--Greengard、Yuan--Lu--Antoine
与 Brown 等论文出发；forward chaining 重点追踪 RtR、Hohage--Soussi 到 Zhang 的
generalized-mode DtN，以及 Barnett--Greengard 到 neighbor/proxy periodization。技术判断
只使用论文原文、正式摘要/全文页或 DOI 信息。bulk band structure、forced scattering、
local defect、fiber/slab 和 line-defect eigenproblem 被严格区分。

证据集包含 25 篇不同的外部学术论文：本地 12 篇与外部新增 13 篇；另有一份内部旧
报告只用于项目溯源，不计入 prior art。20 篇最直接论文进入逐篇笔记和主比较矩阵。

\section{线缺陷 guided-mode 问题的已有理论}

Ammari--Santosa 已建立 line defect 在 photonic band gap 中产生导引谱及色散关系的
数学框架\cite{AmmariSantosa2004}。因此“研究 line-defect guided modes”这一问题层面
显然不是新对象。Fliss 随后针对固定轴向 quasi-momentum，把无界 self-adjoint
eigenproblem 精确化为 defect 邻域上的 nonlinear DtN eigenproblem\cite{Fliss2013}。
其 Theorem 4.5 正是 global/bounded equivalence，Theorem 3.5 又把横向指数衰减率与到
essential spectrum 的距离联系起来。

数值上，Klindworth--Schmidt--Fliss 实现 cell operators、Riccati equation、high-order
FEM、Newton 与 Chebyshev interpolation，并与 supercell 对比；还给出 group velocity
公式\cite{KlindworthSchmidtFliss2014}。Supercell 本身在缺陷模场景有严格的指数收敛结果
\cite{Soussi2005}，但靠近 band edge 时模式衰减弱，所需超胞宽度和建模误差会变差。

因此 draft Theorem 1 的“restriction + gluing”逻辑应被定位为基础 lemma。可能保留的
差异是 full Cauchy relation 不要求 Dirichlet trace 可作为全局坐标，但这必须与已有 RtR
方法比较，并在异常频率给出实质 theorem，而非只说明 graph 可能失效。

\section{周期 half-guide 的 DtN、propagation operator 与 RtR}

Joly--Li--Fliss 在吸收介质中构造 compact propagation operator \(R\)，证明
\(\rho(R)<1\)，并以 stationary Riccati equation 选择物理解\cite{JolyLiFliss2006}。
Coatléven 把该框架用于含 line defect 的强迫问题，允许左右周期外介质不同，并给出
Floquet--Bloch quadrature 与截断误差估计\cite{Coatleven2012}。这两篇都清楚指出：正吸收
给出 coercivity/decay，而 \(\varepsilon\downarrow0\) 的 lossless 情形更困难。

Fliss 2013 的 classical DtN 在一个可数 forbidden set 上不可用。Fliss--Klindworth--
Schmidt 以 Robin-to-Robin map 改写 local problems，使 transparent condition 在这些频率
仍可计算\cite{FlissKlindworthSchmidt2015}。因而“完整 relation 不会在 DtN pole 爆掉”
是一个有潜力的数值设计目标，但不能声称首次绕过 forbidden frequencies。必须回答：

\begin{enumerate}
  \item relation 与 RtR 是否连续/代数等价，还是能覆盖更宽的 threshold/unit-circle 情形；
  \item 离散 relation 的维数、闭性和 inf--sup 常数是否跨 pole 保持；
  \item 与 RtR 相比，条件数、mode count 和 eigenvalue residual 是否更优。
\end{enumerate}

\section{Boundary integral methods 与单元 DtN 的已有工作}

Yuan--Lu--Antoine 已用 transmission BIE 计算任意截面 dielectric cylinders 的 unit-cell
DtN，覆盖两种 scalar polarizations，并用于 bulk band structure\cite{YuanLuAntoine2008}。
所以“用 BIE 代替 FEM”或“只在 interface 上离散”本身不足以构成理论创新。

Barnett--Greengard 的 Theorem 4 在非 exceptional 参数处证明 bulk Bloch eigenvalue
当且仅当 quasiperiodic Müller operator 有非零 kernel\cite{BarnettGreengard2010}。其
Appendix A 通过 swapped-wavenumber complementary transmission problem 证明唯一性，
这与 draft Theorem 2 的证明架构直接重合。该论文也明确用 smallest singular value
扫描 band structure，并进一步以 free-space kernels、相邻 images 和 auxiliary wall
unknowns 绕过 empty-cell/QP Green resonance。

对于一般 transmission BIE，well-posed formulations 和高阶 Nyström 已是成熟工具
\cite{DominguezLyonTurc2016,HaoEtAl2014,Kress1991}。但 Hiptmair--Moiola--Spence 说明：
直接 BIE inverse 可能在 physical solution operator 并未相应放大时出现 spurious
quasi-resonances\cite{HiptmairMoiolaSpence2022}。因此本项目若继续 BIE 路线，不能把
\(\sigma_{\min}\) 的谷值直接解释为 guided mode；必须验证重构场非零、PDE/interface/
transparent residual 同时收敛，并证明或检测 representation nullspace。

\subsection{Theorem 2 的逐式审计}

draft Appendix A 先把总场拆成 homogeneous Rayleigh field 与 outgoing remainder，再
构造 swapped problem。这个总体思路与 Barnett--Greengard Lemmas 9--10 和 Appendix A
高度一致。但在 \texttt{appendixA.tex} 当前源文件第 243 行，第二组 QP potentials 写成
\begin{equation}
  -S_{\rm QP}^{(k)}[v_n^-] + D_{\rm QP}^{(nk)}[v^-].
\end{equation}
若要对内部满足 wavenumber \(k\) 的 \(v^-\) 应用紧邻的 QP interior representation
lemma，第二项必须是 \(D_{\rm QP}^{(k)}[v^-]\)。混用 \(k\) 与 \(nk\) 后，随后声称的
cancellation 不成立。这可能是笔误，但它位于唯一性/存在性证明的关键链条中，必须修正
并重新核对 signs、quasiperiodic phases、Wood anomaly 排除和 auxiliary resonance。

更重要的是，修正笔误并不自动产生新定理：还要精确说明“独立 homogeneous Rayleigh
augmentation”相对于 Barnett 的 range/nullspace theorem 新增了什么，以及为何不会在
interior transmission eigenvalues 或 auxiliary resonances 处丢失唯一性。

\section{Bloch modes、stable subspaces 与 modal completeness}

Joly--Li--Fliss 2006 的 Remark 5.1 曾把 propagation eigenfunction completeness 留作
问题。Hohage--Soussi 2013 的 Theorem 2.2 随后给出 solution space 的 Riesz basis，并
把 translation operator 表示为 infinite Jordan matrix；非平凡 Jordan blocks 至多有限个，
且 Dirichlet/Neumann/mixed traces 继承同一 Jordan 结构\cite{HohageSoussi2013}。

这直接否定 draft Conjecture 1 的普通-eigenvector-only 写法。正确形式不是
\begin{equation}
  \mathcal C_+ = \overline{\operatorname{span}}
  \{\gamma_C u_j:T u_j=\lambda_j u_j,\ |\lambda_j|<1\},
\end{equation}
而应包含 generalized chains \(u_{j,0},\ldots,u_{j,m_j-1}\)，并说明 trace topology、
阈值与 unit-circle modes。Zhang 用 contour/residue 进一步分解 translation operator 和
limiting-absorption solution\cite{Zhang2021}，随后显式构造 generalized-eigenfunction DtN，
证明 modal truncation 指数收敛并给出 FEM 总误差\cite{Zhang2023}。

因此 draft 可有的差异不是“Bloch modes 完备”，而是：从一个高阶 cell scattering
operator 通过 ordered generalized Schur/stable invariant subspace，直接形成 Cauchy
relation，避免显式求一个可能奇异的 Dirichlet block inverse。这个差异必须用连续到离散
的 gap estimate 或 graph-distance estimate 表述，例如
\begin{equation}
  \operatorname{gap}\bigl(\mathcal C_+,
  \mathcal C_{+,N,M}\bigr)
  \le C\bigl(e^{-cM}+E_{\rm Rayleigh}(N)+E_{\rm BIE}(h,p)\bigr),
\end{equation}
而不是只展示若干 eigenvalues。

\section{MFS、quasiperiodic Green function 与 three-cell modification}

Linton 系统整理了 periodic Helmholtz Green function 的 lattice sum、Kummer/Ewald 等
表示\cite{Linton1998}。Luan--Sun--Zhuang 已用 fundamental solutions 加 Rayleigh
expansion 做 quasiperiodic penetrable scattering，并在非 Wood 条件下把场误差控制到
interface residual\cite{LuanSunZhuang2019}。Barnett--Greengard 的相邻 images 思想之后，
Cho--Barnett 明确把 central period 与 immediate left/right neighbors 作为 near field，
proxy sources 表示远处周期 copies，并以 Rayleigh 条件和 Schur complements 消元
\cite{ChoBarnett2015}。

所以 three-cell near-field modification 不是新的算法结构。仍可能有价值的工作是针对
“计算 \(G_{\rm QP}\) 而非直接解 scattering field”的特定 least-squares operator，证明
source radius、proxy number、Rayleigh truncation 与距 Wood anomaly 的联合稳定性。
在没有此类 theorem 前，MFS 应被定位为 enabling implementation。

\section{One-cell scattering matrix 与 Bloch generalized eigenproblem}

bulk bands 通过 transfer/scattering/DtN 的边界 eigenproblem 至少在 2008 年的综述性
论述中已经是标准路线，且 transfer matrix 的指数不稳定早已被指出
\cite{YuanLuAntoine2008}。draft 的 two-port scattering relation
\begin{equation}
  \begin{bmatrix}b_L\\ b_R\end{bmatrix}
  =S(k,\beta)\begin{bmatrix}a_L\\a_R\end{bmatrix},
  \qquad
  \begin{bmatrix}a_R\\b_R\end{bmatrix}
  =\lambda\begin{bmatrix}b_L\\a_L\end{bmatrix}
\end{equation}
经消元成为 generalized eigenproblem，这是标准代数转换。Dohnal--Schweizer 也已经
分析以 Bloch waves 施加周期 outlet radiation condition 的数值方案
\cite{DohnalSchweizer2018}。

可研究的问题是 stable subspace 本身，而非 pencil 形式：在 \(|\lambda|\approx1\)、
multiple/defective multipliers、强 evanescent scale separation 时，如何用 structure-
preserving QZ/ordered Schur 得到可靠 relation，并量化它相对于连续 Cauchy space 的角度。

\section{TE 模式与 divergence-form transmission problems}

TE 模式的自然算子是 divergence form，例如
\begin{equation}
  -\nabla\!\cdot\!\bigl(\varepsilon^{-1}\nabla u\bigr)=k^2u,
\end{equation}
接口条件为 \([u]=0\) 和 \([\varepsilon^{-1}\partial_nu]=0\)。Brown--Hoang--Plum--
Radosz--Wood 已在 line defect、discontinuous coefficients 和负阶 Sobolev/Floquet
框架下证明 arbitrarily weak defects 可向 gap 引入 TE spectrum\cite{BrownEtAl2017}。
Yuan--Lu--Antoine 的 cell DtN/BIE 也包含两种 scalar polarization
\cite{YuanLuAntoine2008}。

所以仅把 TM 的 normal derivative continuity 替换为 weighted flux continuity是 B 级
扩展。它只有作为“统一 TM/TE 的高阶 interface-only solver + 同一透明 relation + 一致
误差/认证框架”的一部分才值得保留；并需明确 principal coefficient jump 对 natural trace、
self-adjoint pairing、Müller weights 与 flux normalization 的影响。

\section{Nonlinear eigenvalue algorithms 与数值检测}

Barnett--Greengard 已用 \(\sigma_{\min}\) 扫描 bulk bands\cite{BarnettGreengard2010}；
Fliss 的数值链已经使用固定点、Newton 与 Chebyshev interpolation
\cite{KlindworthSchmidtFliss2014}。一般 analytic NEP 的 contour/Riesz projection 方法
已经成熟\cite{BinkowskiZschiedrichBurger2020,BrennanEmbreeGugercin2023}，2025 年的
semi-/bi-infinite photonic-crystal 工作还把 contour integral 与 block quasi-Toeplitz/
doubling structure 结合起来\cite{LyuLiLin2025}。

因此 draft 的 grid + recursive dip refinement 只能作为可视化/初值工具。生产求解器应
围绕 contour \(\Gamma\) 计算
\begin{equation}
  P_\Gamma V = \frac{1}{2\pi i}\int_\Gamma
  \mathcal A(z,\beta)^{-1}V\,\mathrm dz,
\end{equation}
并增加至少三层过滤：operator 在 contour 上无 pole；重构物理场非零；PDE/interface/
transparent residual 随离散同步收敛。对 relation-valued block 若非方阵，应先设计稳定的
square linearization 或使用 left/right test spaces，而不是直接对任意 least-squares matrix
套用 contour method。

\section{文献比较矩阵的可读摘要}

完整逐列矩阵见 \texttt{literature\_matrix.csv}。这里按三个问题拆表，避免一张超宽表。

\begin{table}[htbp]
\centering
\caption{直接解决 line-defect guided modes 或透明边界的工作}
\small
\begin{tabularx}{\textwidth}{p{0.22\textwidth} p{0.20\textwidth} p{0.22\textwidth} Y}
\toprule
文献 & 问题 & Exterior treatment & 相对 draft 的结论 \\
\midrule
Ammari--Santosa 2004 & line defect 存在性/色散 & 谱理论 & 问题本身已有数学框架。 \\
Fliss 2013 & line-defect guided eigenmodes & 精确 DtN + Riccati & Theorem 1 和 nonlinear bounded reduction 已覆盖。 \\
Klindworth et al. 2014 & 上述方法数值化 & DtN + high-order FEM & Newton/Chebyshev、group velocity、band-edge 对比均已有。 \\
Fliss et al. 2015 & forbidden frequencies & RtR & 不能仅以“DtN graph 失效”为创新。 \\
Coatléven 2012 & absorbed line-defect scattering & 左右 DtN + FBT & 不同 leads、propagation 与误差估计已有。 \\
\bottomrule
\end{tabularx}
\end{table}

\begin{table}[htbp]
\centering
\caption{BIE、modal 与数值算法先例}
\small
\begin{tabularx}{\textwidth}{p{0.22\textwidth} p{0.25\textwidth} Y}
\toprule
文献 & 已证明/实现 & 对 draft 的约束 \\
\midrule
Barnett--Greengard 2010 & QP Müller nullspace equivalence；swapped problem；\(\sigma_{\min}\) scan & Theorem 2 与扫描高度重合，须明确新 range statement。 \\
Yuan--Lu--Antoine 2008 & BIE 构造 unit-cell DtN，E/H polarization & “FEM 换 BIE”不足。 \\
Hohage--Soussi 2013 & Floquet Riesz basis + Jordan form + trace bases & Conjecture 1 必须包含 generalized modes。 \\
Zhang 2021/2023 & generalized spectrum decomposition；DtN 指数 truncation + FEM error & stable trace truncation 需要达到这一证据标准。 \\
Cho--Barnett 2015 & immediate neighbors + proxies + Rayleigh，Wood robust & three-cell/proxy 架构已有。 \\
Hiptmair et al. 2022 & BIE spurious quasi-resonance & 小奇异值必须配物理残差认证。 \\
Binkowski et al. 2020；Lyu et al. 2025 & contour/Riesz NEP；semi-infinite photonic structure & 网格 scan 不是先进 eigensolver。 \\
\bottomrule
\end{tabularx}
\end{table}

\section{当前 draft 的逐项 prior-art audit}

\begin{longtable}{p{0.29\textwidth} p{0.16\textwidth} p{0.47\textwidth}}
\caption{Claim-by-claim 审计摘要；详细证据见 \texttt{claim\_audit.md}}\\
\toprule
Claim & 状态 & 依据与处理建议 \\
\midrule
\endfirsthead
\toprule
Claim & 状态 & 依据与处理建议 \\
\midrule
\endhead
Theorem 1 global/bounded equivalence & A & Fliss Theorem 4.5 已建立；降为 lemma。\\
DtN graph interpretation & B & 正则处为定义性代数；异常处须与 RtR 比较。\\
Absorbing DtN & A & Joly/Coatléven 已建立，lossless limit 才困难。\\
Rayleigh expansion & A & 非 Wood 情形标准；draft 排除了难点。\\
Theorem 2 exact/unique representation & C & Barnett proof 高度重合；Appendix A 目前有 (k/nk) kernel mismatch。\\
Müller + homogeneous Rayleigh block & B/C & 标准组件的特定耦合；需 Fredholm/kernel-isomorphism theorem。\\
Conjecture 1 completeness & A（需改写） & generalized modes/Jordan chains 的 Riesz basis 已证明；当前写法过强。\\
Conjecture 2 singular iff guided mode & C & 组合性强，且未排除 representation kernel/spurious BIE resonance。\\
MFS QP Green function & B & Luan、Barnett、Cho 等已有；可作基础模块。\\
Scattering-to-Bloch pencil & A/B & 标准转换；稳定 subspace 的连续误差尚可研究。\\
Stable eigentrace relation & C & modal DtN 已有，但与高阶 Müller 的 pole-resistant direct relation coupling 尚未发现完整结果。\\
TE extension & B & divergence form/discontinuous TE theory已有；单纯换权重不新。\\
Smallest singular value scan & A & Barnett 已用；且有漏根/伪根风险。\\
当前 numerical evidence & 预备性 & 一般 line-defect 小节为空，不足以支撑算法 claim。\\
\bottomrule
\end{longtable}

\section{明确已经解决、相近解决与仍不确定的内容}

\subsection{已经解决}

\begin{itemize}
  \item line-defect guided spectrum 的基础数学问题和 gap 中导模存在框架；
  \item global eigenproblem 到 defect neighborhood exact nonlinear DtN problem 的等价；
  \item absorbed half-guide propagation operator、Riccati equation 和 DtN；
  \item 以 RtR 绕过 classical DtN forbidden frequencies；
  \item generalized Floquet modes/Jordan chains 的 Riesz basis 及 trace basis；
  \item generalized modal DtN 的指数截断和一个 FEM 总误差框架；
  \item bulk unit-cell Müller BIE 与 Bloch eigenvalue 的 exact nullspace equivalence；
  \item unit-cell DtN 的 boundary-integral computation；
  \item proxy/neighbor/Rayleigh periodization 与 smallest-singular-value scanning 的基本结构。
\end{itemize}

\subsection{只在不同假设或不同几何下解决}

Müller BIE 的 exact theorem 主要针对 bulk bands；modal completeness/DtN 主要用于 scattering
或抽象 half-guide；contour photonic NEP 的近期工作采用 PEC/semi-infinite block structure；
TE weak-defect work是存在性理论而非界面-only eigensolver。这些不能被错误描述为已经完整
解决 draft 的特定耦合，但足以否定各单一组件的 priority claim。

\subsection{截至本次检索仍不确定}

\begin{itemize}
  \item 一个连续且离散均可认证的 Müller--generalized-Bloch-Cauchy-relation operator，
    在 DtN poles 附近保持稳定并严格排除 representation nullspace；
  \item 针对该特定 operator 的三重误差分解：interface Nyström、Rayleigh truncation、
    Bloch stable-subspace truncation；
  \item Wood anomalies、unit-circle multipliers 与 band-edge thresholds 同时存在时的 relation
    极限；
  \item 相同框架下 dielectric line-defect complex resonances/leaky modes 的 pole tracking。
\end{itemize}

对这些项目只能写：\textbf{截至本次检索尚未发现直接重合工作，但现有证据不足以断言该
方向为首次提出。}

\section{当前 draft 最可能保留的贡献}

最可保留的中心不是 Theorem 1、Conjecture 1、MFS 或 Bloch pencil 单独任何一项，而是
如下交叉命题：

\begin{quote}
在 piecewise-constant dielectric-cylinder line-defect geometry 中，构造 high-order
Müller interface operator 与左右 generalized stable Bloch Cauchy relations 的耦合；证明
离散 kernel 与物理 guided-mode eigenspace 同构，并在 classical DtN forbidden frequencies
附近保持可计算；再用 contour eigensolver 和物理残差排除 representation-induced roots。
\end{quote}

该表述仍需四个最小条件才可成为贡献：修复并重新证明 representation theorem；将
ordinary Bloch modes 改为 generalized invariant subspace；给出至少离散层面的 spurious-
free kernel theorem；完成一般 line-defect convergence/benchmark。少任一项，论文容易被
评价为已有 DtN、Müller、scattering matrix 和 contour method 的串联。

\section{下一步六个候选创新方向}

\subsection{方向一：Spurious-free Müller--Cauchy-relation guided-mode operator}

\begin{description}[style=multiline,leftmargin=4.4cm,labelwidth=4.1cm]
  \item[核心研究问题：] 证明 central-cell Müller block 与左右 outgoing generalized Cauchy relations 的 kernel 与真实 guided eigenspace 一一对应。
  \item[与已有文献的差异：] Fliss 用 FEM--DtN/RtR；Barnett 用 bulk QP Müller；Zhang 用 generalized-mode DtN。目标是三者在 line defect 上的 exact interface-only coupling，并显式排除 representation kernel。
  \item[为什么可能创新：] 本次检索未找到同一几何下的完整 Fredholm/kernel-isomorphism + high-order BIE 结果；Hiptmair 的 spurious 现象使认证不是装饰性细节。
  \item[继续核实的文献：] 更广的 multi-trace/Calderón relation、photonic-crystal-fiber BEM modes、operator-valued NEP 与 pole cancellation。
  \item[主要定理：] Fredholm index zero；物理场重构映射在 kernel 上双射；离散 spectral exactness/no pollution。
  \item[算法：] Kress/Alpert Nyström，ordered generalized Schur stable relation，rank-revealing nullspace，物理 residual filter。
  \item[数值实验：] Fliss geometry、circular/noncircular rods、左右相同/不同 leads、DtN pole 附近；与 FEM--DtN/RtR 独立对比。
  \item[最小可发表单元：] 一个 TM geometry、一个 theorem（kernel 双射）、一套 convergence table、两个 reference dispersion curves。
  \item[难度：] 理论高；实现中高。
  \item[主要风险：] Theorem 2 可能只是一条标准 representation，或 relation block 在异常点不 Fredholm。
  \item[失败时保留：] 经严格验证的高阶 BIE--RtR solver 与 spurious-root diagnostics。
  \item[推荐优先级：] \textbf{1，当前论文主线首选。}
\end{description}

\subsection{方向二：广义稳定 Bloch Cauchy relation 的误差与条件数}

\begin{description}[style=multiline,leftmargin=4.4cm,labelwidth=4.1cm]
  \item[核心研究问题：] 从 discretized one-cell scattering operator 直接逼近 continuous outgoing Cauchy subspace，而不形成 DtN inverse，并给出 gap/angle error。
  \item[差异：] Hohage/Zhang 已给 generalized completeness 和 modal DtN truncation；目标强调 scattering stable invariant subspace、relation coordinates 和跨 DtN pole 条件数。
  \item[创新性依据：] 直接 relation 的条件数优势尚未在检索到的 line-defect BIE 文献中得到证明，但不能据此断言首次。
  \item[继续核实：] wave finite element、symplectic pencils、discrete Riccati invariant subspaces、Krein/flux pairing。
  \item[主要定理：] subspace gap bound；defective multiplier 的 Jordan-chain approximation；远离/接近 unit circle 的条件数界。
  \item[算法：] structure-preserving scattering assembly，QZ/ordered Schur，biorthogonal flux normalization。
  \item[实验：] 人造近重根、\(|\lambda|\approx1\)、DtN Dirichlet block singular；与 DtN/RtR condition number 对比。
  \item[最小单元：] 离散 finite-Rayleigh model 上的定理 + 数值，随后再连接 continuous BIE。
  \item[难度：] 理论很高；实现中。
  \item[风险：] 与 Zhang 2023 的结论等价，只是坐标改变。
  \item[失败时保留：] 稳定的 QZ 模态模块和一套 threshold diagnostics。
  \item[优先级：] \textbf{2；适合作为方向一的理论核心或独立后续。}
\end{description}

\subsection{方向三：Contour NEP + 物理场残差认证}

\begin{description}[style=multiline,leftmargin=4.4cm,labelwidth=4.1cm]
  \item[核心问题：] 用 contour/Riesz projection 找全指定区域的 guided eigenvalues，同时区分 operator poles、representation nullspace 与真实 modes。
  \item[差异：] contour method 本身成熟；差异必须是 relation-valued BIE 的 square realization、pole-safe resolvent 与物理认证。
  \item[创新性：] 可把 draft 最薄弱的 \(\sigma_{\min}\) scan 变成可靠算法，但单纯调用 Beyn/SS 不新。
  \item[继续核实：] meromorphic NEP、BIE pole cancellation、InfBeyn/FEAST、non-square analytic operator pencils。
  \item[主要定理：] contour 内 mode count；无 pole 污染条件；residual filter 对真实 kernel 的一致性。
  \item[算法：] adaptive quadrature、random probing、SVD rank decision、argument-principle/pole check、field reconstruction。
  \item[实验：] multiple/tangent roots、接近 auxiliary resonance、与 dense grid 的漏根/成本对比。
  \item[最小单元：] 固定 TM BIE operator 上的 complete mode count + independently verified eigenpairs。
  \item[难度：] 理论中；实现中。
  \item[风险：] 贡献被视为成熟 eigensolver 的应用。
  \item[失败时保留：] 更可靠的内部研究工具，可立即服务其他方向。
  \item[优先级：] \textbf{1（工程优先），作为方向一的必需算法而非单独主创新。}
\end{description}

\subsection{方向四：Band-edge poorly confined modes 的联合误差控制}

\begin{description}[style=multiline,leftmargin=4.4cm,labelwidth=4.1cm]
  \item[核心问题：] 当 decay rate \(\kappa\to0\) 时，联合控制 BIE、Rayleigh、Bloch-modal 与 eigensolver 误差。
  \item[差异：] Fliss 已有衰减定理与 FEM/supercell 对比；需增加 interface-only relation 的参数一致误差和 adaptive rule。
  \item[创新性：] 若给出 complexity/error law 而非只增加算例，可能有充分算法价值。
  \item[继续核实：] threshold resolvent、slow light/group velocity、supercell asymptotics、spectral pollution。
  \item[主要定理：] 误差常数对 band-edge distance 的显式依赖；adaptive mode count 停止准则。
  \item[算法：] residual-driven Rayleigh/Bloch enrichment，contour branch tracking，high-precision fallback。
  \item[实验：] 从 well-confined 连续追踪到 near-edge，比较 supercell width 与 relation dimension/CPU。
  \item[最小单元：] 一个可重复 benchmark + 三项误差拆分 + 与 Fliss FEM 的独立交叉验证。
  \item[难度：] 理论中高；实现中高。
  \item[风险：] unit-circle threshold 破坏 analytic NEP 假设。
  \item[失败时保留：] 高价值 benchmark 与软件验证数据集。
  \item[优先级：] \textbf{2；数值论文或方向一的强证据。}
\end{description}

\subsection{方向五：左右不同周期 leads 的 interface-only guided-mode solver}

\begin{description}[style=multiline,leftmargin=4.4cm,labelwidth=4.1cm]
  \item[核心问题：] 两侧材料、周期单元或 band gaps 不同，仍计算局域 interface/defect modes。
  \item[差异：] Coatléven 已处理 absorbed scattering，Dohnal--Schweizer 处理不同 periodic outlets；需聚焦 lossless guided eigenmodes 与 BIE relations。
  \item[创新性：] 几何扩展本身不够；必须结合 exact mode existence condition、不同稳定维数匹配和误差认证。
  \item[继续核实：] photonic interface/topological edge states、heterojunction DtN、bi-infinite block methods。
  \item[主要定理：] 两侧 Cauchy relations 交叉的 Fredholm/index 与 mode multiplicity。
  \item[算法：] 独立 left/right cell scattering operators，异维 stable relations，contour eigensolver。
  \item[实验：] identical-lead 回归测试、heterojunction、仅一侧有 gap 的边界案例。
  \item[最小单元：] 两个不同 cylinder lattices 的 TM interface mode + reference FEM。
  \item[难度：] 理论中高；实现中。
  \item[风险：] 被视为 Coatléven/edge-state 方法的直接扩展。
  \item[失败时保留：] 通用双 lead 软件架构。
  \item[优先级：] \textbf{3；适合作为方向一证明后的扩展。}
\end{description}

\subsection{方向六：Wood anomalies、unit-circle multipliers 或 complex leaky resonances}

\begin{description}[style=multiline,leftmargin=4.4cm,labelwidth=4.1cm]
  \item[核心问题：] 越过 draft 当前排除的 Wood/threshold 情形，或把 real guided eigenvalues 延拓到 complex resonances。
  \item[差异：] scattering 中已有用 shifted kernels 或 proxy sources 处理 Wood anomalies 的 formulation；resonance、PML 与 contour 文献也很广，尚未完成与本项目全部交叉的核实。
  \item[创新性：] 难点真实且可能高，但当前证据不足，不能作近期承诺。
  \item[继续核实：] shifted Green functions、threshold radiation conditions、resonance sheets、analytic continuation of stable subspaces。
  \item[主要定理：] threshold relation 极限或 meromorphic continuation；complex pole 与 outgoing field 等价。
  \item[算法：] Wood-robust free-space/proxy BIE，complex contour tracking，sheet-aware branch handling。
  \item[实验：] grazing Rayleigh order、multiplier 穿越 unit circle、由 bound mode 变 leaky mode。
  \item[最小单元：] 一个简化 grating/line-defect threshold model 的理论与算法。
  \item[难度：] 理论极高；实现高。
  \item[风险：] radiation condition 与 analyticity 失败，范围迅速膨胀。
  \item[失败时保留：] Wood-robust QP Green/BIE module。
  \item[优先级：] \textbf{4，探索性后续，不建议当前主线。}
\end{description}

\section{方向排序与不建议作为主创新的内容}

\begin{table}[htbp]
\centering
\caption{候选方向排序（5 为最优/最高；风险 5 为最高）}
\small
\begin{tabularx}{\textwidth}{p{0.30\textwidth} C{1.15cm} C{1.15cm} C{1.15cm} C{1.15cm} C{1.15cm} C{1.35cm} Y}
\toprule
方向 & 新颖潜力 & 可行性 & 时间成本 & 理论风险 & 数值风险 & 代码兼容 & 主线？ \\
\midrule
1 Müller--relation + spurious-free theorem & 4 & 4 & 4 & 4 & 3 & 5 & 是 \\
3 Contour + physical certification & 3 & 5 & 2 & 2 & 2 & 5 & 作为主线算法 \\
2 Stable relation error/conditioning & 4 & 3 & 5 & 5 & 3 & 4 & 可作理论核心 \\
4 Band-edge joint error & 4 & 3 & 4 & 4 & 4 & 4 & 强数值副线 \\
5 Different leads & 3 & 4 & 3 & 3 & 3 & 4 & 后续扩展 \\
6 Wood/complex resonances & 5? & 2 & 5 & 5 & 5 & 3 & 当前否 \\
\bottomrule
\end{tabularx}
\end{table}

不建议作为主创新：仅把 TM 改为 TE；仅重证 Fliss 的 bounded equivalence；只把 FEM 换
成 BIE；只复现 scattering-to-Bloch pencil；只做 MFS/Linton benchmark；只增加新数值图；
只用 contour eigensolver 而不解决 operator pole/spurious kernel。它们可以作为完整方法的
必要组成，却不足以独立支撑 priority claim。

\section{推荐主路线、备用路线与停止条件}

\subsection{推荐主路线}

\textbf{目标：先证明/验证 spurious-free Müller--generalized-Cauchy operator，再实现
contour + physical residual eigensolver，最后以 Fliss/Klindworth geometry 和 band-edge
benchmark 做独立比较。}

\paragraph{第一阶段：两至四周的理论与最小离散审计。}
修正 Theorem 2 的 kernel mismatch；把定理范围缩到清楚的非 Wood、无 auxiliary
resonance 区间；用 Hohage/Zhang 的 generalized modes 重写 Conjecture 1；在固定有限
Rayleigh truncation 上证明 block kernel 重构非零物理场的充要性。代码上增加 field
reconstruction、PDE/interface/port residual 和 null-density/zero-field test。实验先只做
circular rods 的 TM 回归。

\paragraph{第二阶段：数值可靠性与收敛。}
实现 ordered QZ/Schur stable relation 和 contour eigensolver。分别改变 interface nodes、
Rayleigh orders、stable-mode dimension、proxy radius/number，得到三轴收敛表；与一个独立
FEM--DtN/RtR 或充分宽 supercell 对比 eigenvalue、field profile、group velocity。必须包含
一个 near-band-edge poorly confined mode。

\paragraph{第三阶段：论文结构。}
理论部分只保留真正需要的三条结果：representation/kernel-isomorphism、generalized
relation approximation、no-spurious spectral convergence。算法部分给出 contour mode
count 与 residual certification；数值部分给出 standard、DtN-pole-nearby、band-edge 三组。

\paragraph{停止条件。}
若进一步 forward search 找到同一 line-defect Müller--RtR/generalized-relation 的完整
theorem 与同等 high-order numerics，则停止 priority claim，把工作转为复现/软件。若修正
Theorem 2 后仍在普通 non-Wood 参数出现 nonphysical kernel，停止 exact-unique 路线。

\paragraph{转向条件。}
若连续 Fredholm proof 过难，但有限 Rayleigh/BIE block 可严格证明 kernel 双射和 spectral
convergence，则转为“discrete certified solver”较小论文。若 relation 在 DtN pole 无条件数
优势，则直接采用已有 RtR，将贡献收缩为 high-order BIE--RtR + contour certification。

\subsection{较保守备用路线}

实现 high-order Müller BIE--RtR solver（承认 transparent condition 来自已有文献），把
新贡献限定为：可重复的 interface-only implementation、严格的三项离散收敛 study、
contour complete mode count、spurious-root physical residual filter，以及 near-band-edge
与非圆 rods benchmark。该路线理论新颖性较弱，但范围清楚、完成风险低于 general
Cauchy-relation theorem，也与现有代码最兼容。

\section{质量控制与最终结论}

本次审计已区分 scattering 与 eigenvalue、bulk bands 与 line defects、TM 与 TE、DtN/
RtR/trace relation、ordinary 与 generalized Bloch modes。报告没有使用“首次”“全新”或
“此前无人”作为未经证据的结论。外部条目均提供 DOI/URL；本地条目在 BibTeX 中标注
local copy。详细文献元数据、逐篇笔记、搜索记录和每条 draft claim 的证据链均保存在
同一审计目录。

最终判断是：\textbf{当前 draft 的主要组成大多已有先例；现阶段最值得投入的是特定
Müller--\allowbreak generalized-Cauchy coupling 的 spurious-free 认证和可靠 contour eigensolver，
其次是 band-edge 的联合误差/复杂度研究。}在完成 proof repair、generalized-mode 修订和
一般 line-defect 数值验证之前，不建议提交以 Theorems 1--2/Conjectures 1--2 的现有组合
为主创新的论文。

\clearpage
\printbibliography[heading=bibintoc,title={参考文献}]

\end{document}
